% vim:set errorformat="":
\usepackage{amsmath}
\usepackage{amsfonts}
\usepackage{xcolor}
\usepackage{colortbl}
\usepackage{hyperref}
\hypersetup{
    colorlinks=true,
    linkcolor=blue,
    citecolor=blue,
    urlcolor=blue
}
\usepackage{marginnote}
\usepackage{wrapfig}
\usepackage{graphicx}
\usepackage{lipsum}
\usepackage{array}
\usepackage{tabularx}

\usepackage[skins,listings]{tcolorbox}
\usepackage{boolexpr,ifthen}
\definecolor{dkblue}{rgb}{0,0.1,0.5}
\definecolor{lightblue}{rgb}{0,0.5,0.5}
\definecolor{dkgreen}{rgb}{0,0.4,0}
\definecolor{dk2green}{rgb}{0.4,0,0}
\definecolor{dkviolet}{rgb}{0.6,0,0.8}
\definecolor{mantra}{rgb}{1.0,0.5,0}

\def\lstlanguagefiles{defManSSR.tex}
\lstset{language=SSR}
\newcommand{\mcbimpl}[1]{\lstinline/?$_{#1}$/}
\newcommand{\mcbimplm}[1]{\mbox{?}_{#1}}

\newcommand{\C}[1]{\mbox{\lstinline`#1`}}

% Highlights identifiers, |* id *| in underlined red, can we do better?
\lstset{moredelim=[is][\color{red}\bfseries\ttfamily\underbar]{|*}{*|}}

\newcommand\Chapter[2]{\chapter
  [#1\hfil\hbox{}\protect\linebreak{\itshape#2}]
  {#1\\ [2ex] {\Large \itshape #2}}
  \markboth{\MakeUppercase{\chaptername\ \thechapter.\ #1}}{}}

\newcommand{\gotcha}[1]{\framebox[\width]{\parbox{\textwidth}{gotcha: #1}}}

\newcolumntype{V}{>{\centering\arraybackslash} m{0\linewidth} }

\newcommand{\mantra}[1]{
\begin{center}
\begin{tcolorbox}[colframe=mantra,
	center title,
	title={\textcolor{white}{\textsf{\Large\textbf{Mantra}}}}]
\begin{tabular}{lV}
\begin{minipage}{0.86\textwidth}
\vspace{1em}
\large
#1
\end{minipage}
&
\coqhead{1cm}
\end{tabular}
\end{tcolorbox}
\end{center}}

\newcommand{\coqhead}[1]{\includegraphics[width=#1]{../artwork/smallpantade}}

\newcommand{\mcbpplevel}[1]{
\switch[#1=]
\case{1}
\case{2} \hspace{0.6cm} \coqhead{0.6cm}
\case{3} \coqhead{0.6cm} \coqhead{0.6cm}
\otherwise Please use \textbackslash mcbLEVEL
\endswitch
}

\newcommand{\mcbvarreset}{
\mcbLEARN{Please use \textbackslash mcbLEARN}
\mcbREQUIRE{Please use \textbackslash mcbREQUIRE}
\mcbPROVIDE{Please use \textbackslash mcbPROVIDE}
\mcbNOTES{}
\mcbLEVEL{4}}

\newcommand{\mcbLEARN}[1]{\def\mcbvarLEARN{#1}}
\newcommand{\mcbREQUIRE}[1]{\def\mcbvarREQUIRE{#1}}
\newcommand{\mcbPROVIDE}[1]{\def\mcbvarPROVIDE{#1}}
\newcounter{mcbvarLEVEL}
\newcommand{\mcbLEVEL}[1]{\setcounter{mcbvarLEVEL}{#1}}
\newcommand{\mcbNOTES}[1]{\def\mcbvarNOTES{#1}}
\newcounter{mcbvarTMP}

\newcommand{\mcbsection}[1]{
\setcounter{mcbvarTMP}{1}
\ifthenelse{\equal{}{\mcbvarNOTES}}{\setcounter{mcbvarTMP}{0}}{}
\section{#1}
\smash{\raisebox{0.7cm}{\hspace{0.85\textwidth}\mcbpplevel{\value{mcbvarLEVEL}}}}
\marginnote{\small\begin{tabular}{p{3.8cm}}\hline
\textbf{Learns} \mcbvarLEARN \\
\textbf{Requires} \mcbvarREQUIRE \\
\textbf{Provides} \mcbvarPROVIDE \\
\textbf{Level} \arabic{mcbvarLEVEL} \\
\switch[\value{mcbvarTMP}=]
\case{0}
\case{1} \textbf{Note} \mcbvarNOTES \\
\endswitch
\hline
\end{tabular}\\~\\
}\mcbvarreset{}
}

\mcbvarreset{}

\newtcblisting[auto counter]{coq}[1]{
	listing engine=listings,
	listing options={numbers=left,numberstyle=\tiny,
	                 aboveskip=0mm,belowskip=0mm,
		         xrightmargin=0mm,xleftmargin=0mm},
	%on line,
	before={\noindent},
	after={\hfill},
	listing only,
	fonttitle=\sffamily\bfseries,
	#1
}
\newtcblisting{coqout}[1]{
	listing engine=listings,
	listing options={aboveskip=0mm,belowskip=0mm,
	                 xrightmargin=0mm,xleftmargin=0mm},
	%on line,
	before={\noindent},
	after={\hfill},
	listing only,
	fonttitle=\sffamily\bfseries,
	#1
}
